\documentclass[11pt]{article}
\usepackage[margin=1in]{geometry}
\usepackage{amsmath,amssymb}
\usepackage{graphicx}
\usepackage{booktabs}
\usepackage[hidelinks]{hyperref}

\title{Encoder--IMU Fusion for Stance-Leg Odometry\\[2pt]\large A Worked Derivation}
\author{Arda Gencer}
\date{}

\begin{document}
\maketitle

\section{What we actually have, and why neither source is enough alone}

At every control tick we have two independent raw readings: the three joint
encoder angles per leg, $(\theta_{\text{abad}}, \theta_{\text{hip}}, \theta_{\text{knee}})$,
and the IMU's orientation, from which we extract pitch $\theta$ and roll
$\varphi$. Neither one by itself tells us what we actually want, which is the
body's velocity and height in the world frame.

The encoders alone are not enough: forward kinematics on the joint angles
tells us exactly where a foot sits relative to the body's own reference
frame, but it says nothing about which way that frame is tilted relative to
the ground. A leg held at a fixed configuration produces the exact same
joint angles whether the robot is standing level or pitched forward by
30$^\circ$ -- the encoders cannot see the difference.

The IMU alone is not enough either: pitch and roll tell us how the body
frame is tilted, but on their own they say nothing about where the feet
are or whether the robot is translating at all. Orientation and position
are simply different quantities.

So the two sources are complementary in a very literal sense -- one gives a
vector, the other gives the rotation needed to place that vector correctly
in the world. Fusing them here means exactly that: computing a vector from
the encoders, then rotating it with the IMU's angles. Nothing probabilistic
is happening; this is a deterministic geometric composition, not a
Kalman-style statistical fusion.

\section{Step 1 -- encoders give a body-frame vector}

Forward kinematics (the algebraic inverse of the leg's inverse-kinematics
solver) converts one leg's joint angles into the foot's position relative
to that leg's ABAD pivot, in the body's own axes:
\begin{equation}
(f_x, f_y, f_z) = \text{FK}_{3D}(\theta_{\text{abad}}, \theta_{\text{hip}}, \theta_{\text{knee}} \,;\, o_y, s, L_1, L_2, D_{\text{ABAD}})
\end{equation}
where $L_1, L_2$ are the leg's link lengths, $D_{\text{ABAD}}$ is the fixed
ABAD-to-hip arm length, and $o_y, s$ encode which side (left/right) and end
(front/back) of the robot the leg is mounted on. Adding that leg's fixed
mounting offset $(x_0, y_0)$ from the body origin gives the foot position
relative to the body origin itself,
\begin{equation}
\mathbf{p}_{\text{foot}}^{\,B} = (f_x + x_0,\; f_y + y_0,\; f_z).
\end{equation}
The superscript $B$ is deliberate: this vector is only correct when read in
body-frame axes. We do not yet know how those axes sit relative to the
world.

\section{Step 2 -- the IMU supplies the missing rotation}

All three orientation angles -- pitch $\theta$, roll $\varphi$, and yaw
$\psi$ -- come straight from the IMU's orientation quaternion. All three are
needed: leaving any one out means the rotation only correctly places a
vector for whichever value of that angle happens to be zero, and silently
misplaces it otherwise. Section~\ref{sec:yaw-fix} says more about where
these three angles come from and what using them here does and does not
assume. The body-to-world rotation used is
\begin{equation}
R(\psi, \theta, \varphi) = R_z(\psi)\,R_y(\theta)\, R_x(\varphi) =
\footnotesize
\begin{pmatrix}
\cos\psi\cos\theta & \sin\varphi\sin\theta\cos\psi - \sin\psi\cos\varphi & \sin\varphi\sin\psi + \sin\theta\cos\varphi\cos\psi \\
\sin\psi\cos\theta & \sin\varphi\sin\psi\sin\theta + \cos\varphi\cos\psi & \sin\psi\sin\theta\cos\varphi - \sin\varphi\cos\psi \\
-\sin\theta & \sin\varphi\cos\theta & \cos\varphi\cos\theta
\end{pmatrix}
\normalsize,
\label{eq:R-full}
\end{equation}
with $\theta$ = pitch, $\varphi$ = roll, $\psi$ = yaw (composed as roll,
then pitch, then yaw -- matrix composition applies right to left, so a
body-frame vector is rotated by roll first, then pitch, then yaw). Applying
it to the encoder-derived vector re-expresses the same physical offset in
world-frame axes:
\begin{equation}
\mathbf{p}_{\text{foot}}^{\,W} = R(\psi,\theta,\varphi)\,\mathbf{p}_{\text{foot}}^{\,B}
= (w_x, w_y, w_z).
\end{equation}
This single matrix multiplication is the entire fusion step: an
encoder-derived vector, rotated by an IMU-derived matrix.

\section{Where the orientation angles come from, and their limits}
\label{sec:yaw-fix}

It is worth being precise about what $(\psi,\theta,\varphi)$ actually are
in this implementation, since it affects how the whole estimator should be
read. All three are taken directly from the IMU's orientation quaternion --
that is, from the simulator's exact ground-truth attitude, not from a value
derived through any sensor-fusion process. Using the full rotation is
still the mathematically correct thing to do with that reading, regardless
of where it came from: a rotation matrix missing any one of the three
angles only places a vector correctly when that particular angle happens
to be zero, and misplaces it -- specifically, misdirects horizontal
velocity and position, since height is unaffected by yaw entirely
(Section~\ref{sec:height}) -- by however much it actually isn't.

Whether this orientation is realistically obtainable is a separate
question worth answering honestly. On real hardware, roll and pitch are
reasonably recoverable on their own, since gravity gives a fixed, known
reference direction that a stationary or slowly-moving accelerometer can
sense. Yaw is a different matter: Bloesch et al. (\emph{State Estimation
for Legged Robots -- Consistent Fusion of Leg Kinematics and IMU}, RSS
2012) show that yaw and absolute position together form a manifold that is
\emph{provably unobservable} from IMU and leg kinematics alone -- only
roll, pitch, and body velocity can be recovered from that sensor set.
Recovering yaw for real would require an additional absolute reference this
project does not have, such as a magnetometer (workable in principle, but
vulnerable in practice to interference from the leg actuators' own motor
currents) or a vision/LIDAR-based heading estimate. So this estimator
should be read precisely for what it is: a correct treatment of a known
orientation, where that orientation itself remains a simulation privilege
rather than something this sensor suite could produce on a real robot.

\section{Height falls out of a flat-ground constraint}
\label{sec:height}

If the ground is flat at world $z=0$ and the current foot is planted on it,
then the body's height plus the foot's world-frame $z$-offset from the body
must sum to zero:
\begin{equation}
z_{\text{body}} + w_z = 0 \quad\Longrightarrow\quad z_{\text{body}} = -w_z.
\end{equation}
Notice this is a single algebraic read-out of Step 2's result -- nothing is
integrated, so this estimate carries no accumulating drift; each tick is a
fresh, independent computation. It is also, conveniently, completely
insensitive to yaw: a rotation about the vertical axis only ever mixes the
$x$ and $y$ components of a vector into each other, it can never move
anything into or out of $z$. So of the three orientation angles discussed
in Section~\ref{sec:yaw-fix}, only pitch and roll matter here -- yaw affects
horizontal velocity and position only.

\section{Velocity comes from differentiating a constancy constraint}

The physical statement that gives us velocity is: while a foot is planted
and not slipping, its world position does not move. Writing the body's
world position as $\mathbf{p}_{\text{body}}^{\,W}(t)$,
\begin{equation}
\mathbf{p}_{\text{body}}^{\,W}(t) + R(t)\,\mathbf{p}_{\text{foot}}^{\,B}(t) = \text{constant},
\end{equation}
where $R(t)$ is shorthand for the full rotation of Eq.~\eqref{eq:R-full},
evaluated at whatever $(\psi,\theta,\varphi)$ the IMU reports at time $t$.
Differentiating both sides with respect to time,
\begin{equation}
\mathbf{v}_{\text{body}}^{\,W}(t) = -\frac{d}{dt}\Big[R(t)\,\mathbf{p}_{\text{foot}}^{\,B}(t)\Big]
= -\Big[\dot{R}(t)\,\mathbf{p}_{\text{foot}}^{\,B}(t) + R(t)\,\dot{\mathbf{p}}_{\text{foot}}^{\,B}(t)\Big],
\label{eq:product-rule}
\end{equation}
by the product rule. There are genuinely two contributions here: one from
the leg's own joint angles changing ($\dot{\mathbf{p}}_{\text{foot}}^{\,B}$), and
one from the \emph{body itself rotating} while the foot vector is held
($\dot R$). The first term is approximated with a finite difference of the
body-frame foot position, exactly as the geometry suggests. The second term
is not dropped: using the standard rigid-body identity
$\dot R(t) = R(t)\,[\boldsymbol\omega^{b}(t)]_\times$, where $\boldsymbol\omega^{b}(t)$
is the body-frame angular velocity the gyro reports directly (the same
IMU message already supplies orientation and angular velocity together, so
no separate estimation step is needed), this term can be rewritten as
\begin{equation}
\dot R(t)\,\mathbf{p}_{\text{foot}}^{\,B}(t) = R(t)\Big(\boldsymbol\omega^{b}(t) \times \mathbf{p}_{\text{foot}}^{\,B}(t)\Big),
\end{equation}
a cross product computed entirely in body frame, which folds directly into
the same finite difference before the single rotation is applied:
\begin{equation}
\mathbf{v}_{\text{body}}^{\,W} \;\approx\; -R_k \left[\frac{\mathbf{p}_{\text{foot}}^{\,B}(t_k) - \mathbf{p}_{\text{foot}}^{\,B}(t_{k-1})}{\Delta t} \;+\; \boldsymbol\omega^{b}(t_k) \times \mathbf{p}_{\text{foot}}^{\,B}(t_k)\right].
\label{eq:vel-approx}
\end{equation}
This is the only approximation left: $\boldsymbol\omega^{b}$ is evaluated once,
at the current tick, rather than integrated exactly across
$[t_{k-1}, t_k]$, so a small amount of $\Delta t$-discretization error
remains -- but the rotation-transport term itself is no longer thrown away.
A synthetic check confirms this matters in practice: for a body translating
at a known velocity while also rotating at a realistic
$0.1$--$0.6\,\text{rad/s}$, dropping this term (as an earlier version of
this estimator did) produced a velocity error of about
$0.17\,\text{m/s}$ on a true speed of $0.37\,\text{m/s}$; including it
reduced that error to about $0.0014\,\text{m/s}$, consistent with the
residual being pure discretization error rather than a remaining
systematic bias.

Once the bracketed term $\mathbf{v}_{\text{rel}}^{\,W} = R_k\big[(\mathbf{p}_k - \mathbf{p}_{k-1})/\Delta t + \boldsymbol\omega^b(t_k)\times\mathbf{p}_k\big]$ is
computed, the stance-foot constraint flips its sign: the foot
\emph{appearing} to move by $\mathbf{v}_{\text{rel}}^{\,W}$ relative to the body
means the body itself is moving by exactly the opposite,
\begin{equation}
\mathbf{v}_{\text{body}}^{\,W} = -\mathbf{v}_{\text{rel}}^{\,W}.
\end{equation}

\section{Worked example}
\label{sec:worked-example}

Take the front-left leg ($o_y = +1$, $s = +1$), with
$L_1 = L_2 = 0.2\,\text{m}$, $D_{\text{ABAD}} = 0.1\,\text{m}$, and mounting
offset $(x_0, y_0) = (0.15,\, 0.113)\,\text{m}$. Suppose two consecutive
ticks, $20\,\text{ms}$ apart, read:

\begin{center}
\begin{tabular}{@{}lccc@{}}
\toprule
 & $\theta_{\text{abad}}$ & $\theta_{\text{hip}}$ & $\theta_{\text{knee}}$ \\
\midrule
tick $k-1$ & $0.050$ & $-0.100$ & $0.300$ \\
tick $k$   & $0.050$ & $-0.080$ & $0.290$ \\
\bottomrule
\end{tabular}
\qquad
\begin{tabular}{@{}lccc@{}}
\toprule
 & pitch $\theta$ & roll $\varphi$ & yaw $\psi$ \\
\midrule
tick $k-1$ & $0.040$ & $0.015$ & $0.350$ \\
tick $k$   & $0.050$ & $0.020$ & $0.360$ \\
\bottomrule
\end{tabular}
\end{center}
The robot has turned about $20^\circ$ since the estimator started ($\psi
\approx 0.35$--$0.36\,\text{rad}$), so this example also shows why the full
rotation of Eq.~\eqref{eq:R-full} genuinely needs all three angles rather
than just pitch and roll.

\noindent \textbf{Step 1.} Forward kinematics at each tick gives the
foot position relative to the body origin, in body-frame axes -- this step
never involves the IMU, so it is unchanged from before:
\begin{align}
\mathbf{p}_{\text{foot}}^{\,B}(t_{k-1}) &= (0.169767,\ 0.132618,\ -0.389523)\ \text{m}, \\
\mathbf{p}_{\text{foot}}^{\,B}(t_{k})   &= (0.175709,\ 0.132615,\ -0.389475)\ \text{m}.
\end{align}
Since the ABAD angle barely changed between ticks (both $0.05\,\text{rad}$)
while hip and knee opened up slightly, almost all of the change shows up in
$f_x$ -- exactly what we'd expect: swinging the hip/knee in their own plane
moves the foot mostly along the leg's forward direction, not sideways.

\noindent \textbf{Step 2.} Rotating each by that tick's own full
$R(\psi,\theta,\varphi)$ (Eq.~\eqref{eq:R-full}):
\begin{align}
\mathbf{p}_{\text{foot}}^{\,W}(t_{k-1}) &= (0.097319,\ 0.182905,\ -0.393968)\ \text{m}
&\Rightarrow\ z_{\text{body}} &= 0.393968\ \text{m}, \\
\mathbf{p}_{\text{foot}}^{\,W}(t_{k})   &= (0.096698,\ 0.186390,\ -0.395043)\ \text{m}
&\Rightarrow\ z_{\text{body}} &= 0.395043\ \text{m}.
\end{align}
Both height values agree with each other to well within a millimeter
despite the body having pitched and rolled slightly between ticks, which is
exactly the point of rotating before reading off $z$ rather than reading
$f_z$ directly. Consistent with Section~\ref{sec:yaw-fix}'s point about
yaw and height, recomputing this same example at $\psi = 0$ instead leaves
$z_{\text{body}}$ completely unchanged at both ticks -- only $w_x$ and
$w_y$ move, since a $20^\circ$ heading change substantially redistributes a
vector between $x$ and $y$.

\noindent \textbf{Step 3 -- velocity, before the rotation correction.} The
body-frame finite difference is unchanged from Step 1 (it doesn't involve
the IMU either):
\begin{equation}
\frac{\mathbf{p}_{\text{foot}}^{\,B}(t_k) - \mathbf{p}_{\text{foot}}^{\,B}(t_{k-1})}{\Delta t}
= (0.297093,\ -0.000119,\ 0.002378)\ \text{m/s}.
\end{equation}

\noindent \textbf{Step 4 -- the rotation-transport correction.} This
example's pitch, roll, and yaw all change noticeably from one tick to the
next, so it is also a useful check of $\boldsymbol\omega^{b}\times
\mathbf{p}_{\text{foot}}^{\,B}$. In place of a real gyro sample, approximate
$\boldsymbol\omega^{b}$ here by finite-differencing the tick-to-tick angles
themselves (a reasonable stand-in at these small angles, though not what
the actual implementation does -- it reads $\boldsymbol\omega^{b}$ directly
from the IMU message):
\begin{equation}
\boldsymbol\omega^{b} \approx \Big(\tfrac{0.020-0.015}{0.02},\ \tfrac{0.050-0.040}{0.02},\ \tfrac{0.360-0.350}{0.02}\Big)
= (0.25,\ 0.50,\ 0.50)\ \text{rad/s}.
\end{equation}
Crossing this with $\mathbf{p}_{\text{foot}}^{\,B}(t_k)$ gives a correction
of $(-0.261045,\ 0.185223,\ -0.054701)\ \text{m/s}$ -- \emph{larger} than the
finite-difference term itself, and enough to change which direction the
combined bracket points:
\begin{equation}
\frac{\mathbf{p}_{\text{foot}}^{\,B}(t_k) - \mathbf{p}_{\text{foot}}^{\,B}(t_{k-1})}{\Delta t} + \boldsymbol\omega^{b}(t_k)\times\mathbf{p}_{\text{foot}}^{\,B}(t_k)
= (0.036055,\ 0.185073,\ -0.052301)\ \text{m/s}.
\end{equation}
Rotating this corrected bracket by the current tick's full
$R(\psi_k,\theta_k,\varphi_k)$ and flipping sign gives the corrected
estimate,
\begin{equation}
\mathbf{v}_{\text{body}}^{\,W} = (0.034123,\ -0.185984,\ 0.050330)\ \text{m/s},
\end{equation}
compared with $(-0.277871,\ -0.104413,\ 0.012476)\ \text{m/s}$ from ignoring
the rotation term entirely -- not just a small correction here, but one
that changes the sign of two of the three components. The reason this
particular example is so sensitive is that its angles were originally
chosen to change quickly from tick to tick purely to make yaw's effect on
the rotation matrix visible (Step 2); a real gait tick with a slower,
steadier body rotation would generally see a smaller correction than this,
though how small is an empirical question for the actual logged
$\boldsymbol\omega^{b}$, not something to assume.

\section{Combining legs, and the final integration}

When more than one foot is in stance simultaneously, each produces its own
independent $(z_{\text{body}}, \mathbf{v}_{\text{body}}^{\,W})$ estimate via the
steps above. These are combined with a plain, unweighted arithmetic mean.
It is worth being honest that this is a simplification: a leg near full
extension is geometrically more sensitive to a given joint-angle error than
one near mid-range, so a statistically rigorous fusion would weight each
leg's contribution by its estimated reliability (e.g.\ a weighted
least-squares or Kalman update) rather than trusting every stance leg
equally.

Finally, horizontal position is obtained with a single Euler integration of
the fused velocity,
\begin{equation}
x_{\text{body}}(t_k) = x_{\text{body}}(t_{k-1}) + v_{x}\,\Delta t,
\end{equation}
and likewise for $y$. This is the only integration anywhere in the
pipeline -- height never accumulates error since it is read fresh each
tick, and velocity itself is obtained by differentiation rather than
integration, so horizontal position is one integration away from a
non-drifting quantity rather than two integrations away from a
double-drifting one, as a raw accelerometer-based estimate would be.

\end{document}
